% Options for packages loaded elsewhere
\PassOptionsToPackage{unicode}{hyperref}
\PassOptionsToPackage{hyphens}{url}
%
\documentclass[
]{article}
\usepackage{amsmath,amssymb}
\usepackage{iftex}
\ifPDFTeX
  \usepackage[T1]{fontenc}
  \usepackage[utf8]{inputenc}
  \usepackage{textcomp} % provide euro and other symbols
\else % if luatex or xetex
  \usepackage{unicode-math} % this also loads fontspec
  \defaultfontfeatures{Scale=MatchLowercase}
  \defaultfontfeatures[\rmfamily]{Ligatures=TeX,Scale=1}
\fi
\usepackage{lmodern}
\ifPDFTeX\else
  % xetex/luatex font selection
\fi
% Use upquote if available, for straight quotes in verbatim environments
\IfFileExists{upquote.sty}{\usepackage{upquote}}{}
\IfFileExists{microtype.sty}{% use microtype if available
  \usepackage[]{microtype}
  \UseMicrotypeSet[protrusion]{basicmath} % disable protrusion for tt fonts
}{}
\makeatletter
\@ifundefined{KOMAClassName}{% if non-KOMA class
  \IfFileExists{parskip.sty}{%
    \usepackage{parskip}
  }{% else
    \setlength{\parindent}{0pt}
    \setlength{\parskip}{6pt plus 2pt minus 1pt}}
}{% if KOMA class
  \KOMAoptions{parskip=half}}
\makeatother
\usepackage{xcolor}
\setlength{\emergencystretch}{3em} % prevent overfull lines
\providecommand{\tightlist}{%
  \setlength{\itemsep}{0pt}\setlength{\parskip}{0pt}}
\setcounter{secnumdepth}{-\maxdimen} % remove section numbering
\ifLuaTeX
  \usepackage{selnolig}  % disable illegal ligatures
\fi
\usepackage{bookmark}
\IfFileExists{xurl.sty}{\usepackage{xurl}}{} % add URL line breaks if available
\urlstyle{same}
\hypersetup{
  hidelinks,
  pdfcreator={LaTeX via pandoc}}

\author{}
\date{}

\begin{document}

\section{Semantic Obstruction and Irreducible
Residue}\label{semantic-obstruction-and-irreducible-residue}

\subsection{A Category-Theoretic Framework for Representation-Invariant
Constraints, Traversability, and Local-to-Global
Failure}\label{a-category-theoretic-framework-for-representation-invariant-constraints-traversability-and-local-to-global-failure}

\subsubsection{Author}\label{author}

Lu Semita (and other Authors to be named later upon formalizing their
input)

\subsubsection{Prefatory Note}\label{prefatory-note}

This manuscript is a work in progress and is being published strictly
for the limited purpose of demonstrating the existence of this work as
of this date, in order to help prevent competitive claims of prior art.
It is offered as a timestamped record of independent work and
development, and not as a claim of exclusive ownership.

\subsubsection{Abstract}\label{abstract}

This paper introduces a general framework for analyzing constraints
relative to representation. A semantic obstruction is defined as
obstruction measured within a chosen representational scheme, while the
irreducible residue is the orbit-minimized obstruction over all
admissible re-descriptions. This yields a canonical decomposition of
obstruction into representation-dependent and representation-invariant
components, establishes orbit invariance, proves existence under mild
compactness and lower-semicontinuity hypotheses, and produces a
trichotomy classifying system-target pairs as removable, irreducible and
fatal, or irreducible and traversable. The framework is formulated
categorically in terms of a representation groupoid, then lifted to
gauge-theoretic, sheaf-theoretic, and higher-categorical settings.
Worked examples are given in gauge theory, nonholonomic control, sheaf
gluing, and optimization. A final appendix shows how the framework can
serve as a universal front-end for broad mathematical research programs,
including but not limited to work on Yang-Mills, Navier-Stokes, the
Riemann Hypothesis, P versus NP, and structurally similar problems. The
framework is general and does not itself solve any particular open
problem; rather, it formalizes a reusable invariant structure for
distinguishing removable obstruction from irreducible but potentially
traversable residue.

\subsubsection{Keywords}\label{keywords}

semantic obstruction, irreducible residue, obstruction theory,
representation groupoid, gauge invariance, cocycle obstruction,
traversability, local-to-global failure, category theory, higher
category theory, nonholonomic control, sheaf theory

\subsubsection{2020 Mathematics Subject
Classification}\label{mathematics-subject-classification}

Primary: 18B40, 18N99, 55S35, 58B25

Secondary: 55R65, 58A14, 93B29, 68Q17, 81T13

Suggested alternatives depending on venue and emphasis: - 18A05 (general
category theory) - 18F20 (categories and geometry) - 55R10 (fiber
bundles) - 55N30 (sheaf cohomology) - 58Z05 (applications to physics) -
93C10 (nonlinear systems and control) - 81T25 (quantized field theory on
curved backgrounds or topological aspects)

\begin{center}\rule{0.5\linewidth}{0.5pt}\end{center}

\section{1. Introduction}\label{introduction}

Many hard problems arise through an interplay between structure and
description. A system may appear blocked because it has been placed in a
poor coordinate system, an inconvenient encoding, or an unsuitable local
model. In that case, the obstruction is not intrinsic to the system
alone; it is relative to the representation. On the other hand, some
obstruction remains no matter how the system is re-described within the
admissible class. That persistent remainder is the real structural
content of the problem.

The aim of this paper is to formalize that distinction.

The central thesis is simple:

The mathematically meaningful part of an obstruction is not its value in
any one representation, but the minimum obstruction that survives across
all admissible re-descriptions.

This minimum is called the irreducible residue.

Once this is formalized, a new trichotomy becomes visible. A
system-target pair may be:

\begin{enumerate}
\def\labelenumi{\arabic{enumi}.}
\tightlist
\item
  removable, meaning the obstruction can be driven to zero by admissible
  re-description;
\item
  irreducible and fatal, meaning a positive residue remains and no
  successful admissible traversal exists;
\item
  irreducible and traversable, meaning a positive residue remains but
  the target can still be reached by an admissible representation-action
  pair.
\end{enumerate}

The third regime is especially important. In practice, many systems are
not made successful by removing all obstruction, but by aligning with
what cannot be removed. This phenomenon appears in control, topology,
gauge theory, engineering, and computation. The present framework is
built to formalize precisely that regime.

The paper proceeds as follows. Section 2 defines the representation
groupoid and obstruction functional. Section 3 introduces the
irreducible residue and apparent obstruction. Section 4 proves the first
structural results: decomposition, invariance, monotonicity, and
existence. Section 5 introduces traversability and proves the
trichotomy. Section 6 develops the categorical formulation. Section 7
specializes to cocycle and sheaf-theoretic obstruction. Section 8 gives
a gauge-theoretic formulation. Section 9 lifts the theory to
higher-categorical language. Section 10 gives worked examples with
proofs. Section 11 discusses design principles and interpretation.
Section 12 concludes. Appendix A provides an implementation guide for
using the framework as a universal front-end to major research programs
and analogous problem families.

\section{2. Basic Setup}\label{basic-setup}

\subsection{2.1. System states, targets, and
representations}\label{system-states-targets-and-representations}

Let \(X\) be a set, topological space, manifold, moduli space, or other
admissible class of system states.

Let \(T\) denote a target. The target may be: - a property to be
achieved, - a simplification to be realized, - a traversal objective, -
a consistency condition, - a desired construction, - a regularity
condition, - a reduction in complexity, - or a spectral, topological, or
dynamical requirement.

For each \(x \in X\), let \(\mathcal{R}(x)\) be a nonempty collection of
representations of \(x\). A representation may be: - a coordinate chart,
- a local trivialization, - an encoding, - a decomposition, - a model, -
a computational description, - a gauge choice, - a coarse-graining, - or
any other admissible mode of description.

We write \(r \in \mathcal{R}(x)\) for a representation of \(x\).

\subsection{2.2. The representation
groupoid}\label{the-representation-groupoid}

We assume the admissible changes of representation are organized into a
groupoid \(\mathcal{G}\), called the representation groupoid.

For each fixed \(x \in X\), the representations \(\mathcal{R}(x)\) form
a connected component, or orbit, in \(\mathcal{G}\). Thus: - objects of
\(\mathcal{G}\) are representations, - morphisms are admissible
re-descriptions, - every morphism is invertible.

We say \(r_1 \sim r_2\) if \(r_1\) and \(r_2\) lie in the same orbit of
\(\mathcal{G}\).

The point of using a groupoid rather than a group is that different
states may admit different representation classes, and not every
representation needs to be globally comparable to every other one.

\subsection{2.3. Obstruction functional}\label{obstruction-functional}

Fix a target \(T\).

An obstruction functional is a map \[
O: \{(r,x) : x \in X,\ r \in \mathcal{R}(x)\} \to [0,\infty]
\] which we write as \[
O(r,x;T).
\]

Interpretationally, \(O(r,x;T)\) measures how obstructed the target
\(T\) is when the system \(x\) is viewed through the representation
\(r\).

Depending on the domain, \(O\) may measure: - failure of commutativity,
- cocycle defect, - nontrivial holonomy, - curvature residue, -
nontrivial topological class, - cost of control, - optimization
infeasibility, - descriptive complexity, - symmetry mismatch, -
concentration defect, - unresolved singularity signal, - lower-bound
witness, - or another obstruction quantity.

No single interpretation is built into the definition.

\section{3. Irreducible Residue and Apparent
Obstruction}\label{irreducible-residue-and-apparent-obstruction}

\subsection{3.1. Irreducible residue}\label{irreducible-residue}

\subsubsection{Definition 3.1}\label{definition-3.1}

Let \(x \in X\). The irreducible residue of \(x\) relative to target
\(T\) is \[
\boxed{
\mathcal{I}(x;T) := \inf_{r \in \mathcal{R}(x)} O(r,x;T)
}
\]

This is the minimal obstruction remaining after optimizing over all
admissible representations of \(x\).

If \(\mathcal{I}(x;T) = 0\), then the obstruction is removable at the
representational level.

If \(\mathcal{I}(x;T) > 0\), then some positive obstruction survives
every admissible re-description.

\subsection{3.2. Apparent obstruction}\label{apparent-obstruction}

\subsubsection{Definition 3.2}\label{definition-3.2}

For \(r \in \mathcal{R}(x)\) such that \(O(r,x;T) < \infty\), define the
apparent obstruction by \[
\mathcal{A}(r,x;T) := O(r,x;T) - \mathcal{I}(x;T).
\]

This is the portion of the obstruction visible in representation \(r\)
that is not forced by the orbit-minimized residue.

The decomposition \[
O(r,x;T) = \mathcal{A}(r,x;T) + \mathcal{I}(x;T)
\] is the basic algebraic split between representational artifact and
structural remainder.

\section{4. Main Structural Theorems}\label{main-structural-theorems}

\subsection{4.1. Canonical decomposition}\label{canonical-decomposition}

\subsubsection{Theorem 4.1}\label{theorem-4.1}

For every \(x \in X\), target \(T\), and representation
\(r \in \mathcal{R}(x)\) with finite obstruction, \[
O(r,x;T) = \mathcal{A}(r,x;T) + \mathcal{I}(x;T),
\] with \[
\mathcal{A}(r,x;T) \ge 0.
\]

\subsubsection{Proof}\label{proof}

By definition of infimum, \[
\mathcal{I}(x;T) \le O(r,x;T)
\] for every \(r \in \mathcal{R}(x)\). Therefore \[
O(r,x;T) - \mathcal{I}(x;T) \ge 0.
\] By Definition 3.2, this difference is precisely
\(\mathcal{A}(r,x;T)\). Rearranging yields the decomposition. QED.

\subsection{4.2. Orbit invariance}\label{orbit-invariance}

\subsubsection{Theorem 4.2}\label{theorem-4.2}

The irreducible residue \(\mathcal{I}(x;T)\) is invariant under
admissible re-description. Equivalently, it depends only on the orbit
\(\mathcal{R}(x) \subseteq \mathcal{G}\).

\subsubsection{Proof}\label{proof-1}

The definition of \(\mathcal{I}(x;T)\) takes the infimum over all
representations in the orbit \(\mathcal{R}(x)\). Any admissible
transformation within that orbit preserves the set over which the
infimum is taken. Therefore the value of the infimum is unchanged. QED.

\subsection{4.3. Monotonicity under restriction of admissible
class}\label{monotonicity-under-restriction-of-admissible-class}

\subsubsection{Theorem 4.3}\label{theorem-4.3}

Let \(\mathcal{R}_1(x)\) and \(\mathcal{R}_2(x)\) be two admissible
representation classes for the same state \(x\), with \[
\mathcal{R}_2(x) \subseteq \mathcal{R}_1(x).
\] Then \[
\mathcal{I}_{\mathcal{R}_2}(x;T) \ge \mathcal{I}_{\mathcal{R}_1}(x;T).
\]

\subsubsection{Proof}\label{proof-2}

The infimum over a smaller set cannot be less than the infimum over a
larger set. QED.

\subsection{4.4. Existence of minimizing
representation}\label{existence-of-minimizing-representation}

\subsubsection{Theorem 4.4}\label{theorem-4.4}

Assume: 1. \(\mathcal{R}(x)\) is compact, 2. \(r \mapsto O(r,x;T)\) is
lower semicontinuous.

Then there exists \(r^\ast \in \mathcal{R}(x)\) such that \[
O(r^\ast,x;T) = \mathcal{I}(x;T).
\]

\subsubsection{Proof}\label{proof-3}

A lower semicontinuous function on a compact space attains its infimum.
QED.

\subsubsection{Remark}\label{remark}

This theorem supplies a standard sufficient condition for the
irreducible residue to be realized by an actual representation rather
than only as a limiting value.

\section{5. Traversability and the
Trichotomy}\label{traversability-and-the-trichotomy}

\subsection{5.1. Success predicate and admissible
actions}\label{success-predicate-and-admissible-actions}

Let \(\mathcal{U}(x)\) be a set of admissible actions, controls, or
procedures available for state \(x\).

Let \[
P(r,u,x;T)
\] be a success predicate taking values in
\(\{\text{true},\text{false}\}\), or more generally a performance
functional whose threshold determines whether the target has been
achieved.

\subsection{5.2. Traversability}\label{traversability}

\subsubsection{Definition 5.1}\label{definition-5.1}

We say that \((x,T)\) is traversable if there exist \[
r \in \mathcal{R}(x), \quad u \in \mathcal{U}(x)
\] such that \[
P(r,u,x;T)
\] holds.

The key point is that traversability is logically independent of whether
\(\mathcal{I}(x;T)\) vanishes.

\subsection{5.3. Trichotomy}\label{trichotomy}

\subsubsection{Theorem 5.2}\label{theorem-5.2}

For every system-target pair \((x,T)\), exactly one of the following
holds:

\begin{enumerate}
\def\labelenumi{\arabic{enumi}.}
\item
  Removable: \[
  \mathcal{I}(x;T)=0.
  \]
\item
  Irreducible and fatal: \[
  \mathcal{I}(x;T)>0
  \quad\text{and}\quad
  (x,T)\text{ is not traversable.}
  \]
\item
  Irreducible and traversable: \[
  \mathcal{I}(x;T)>0
  \quad\text{and}\quad
  (x,T)\text{ is traversable.}
  \]
\end{enumerate}

\subsubsection{Proof}\label{proof-4}

Either \(\mathcal{I}(x;T)=0\) or \(\mathcal{I}(x;T)>0\). In the first
case, we are in case (1). In the second case, either there exists an
admissible successful pair \((r,u)\), or none exists. These alternatives
are mutually exclusive and exhaustive, giving cases (2) and (3). QED.

\subsubsection{Interpretation}\label{interpretation}

This theorem formalizes the third regime that is frequently missing in
standard treatments: non-removable obstruction need not imply
impossibility.

\section{6. Category-Theoretic
Formulation}\label{category-theoretic-formulation}

\subsection{6.1. Obstruction as an order-valued
functor}\label{obstruction-as-an-order-valued-functor}

Let \(([0,\infty], \le)\) be viewed as a thin category.

Then an obstruction assignment may be regarded as a functor-like
structure \[
O_T: \mathcal{R}(x) \to ([0,\infty],\le)
\] provided one imposes a compatibility condition on morphisms. In the
most basic form, the functor sends each object \(r\) to the value
\(O(r,x;T)\). If morphisms preserve or monotonize obstruction values,
then \(O_T\) behaves as a genuine functor into an ordered target
category.

The irreducible residue is then the lower envelope of this order-valued
assignment on the connected component \(\mathcal{R}(x)\).

\subsection{6.2. Groupoid viewpoint}\label{groupoid-viewpoint}

The groupoid viewpoint is especially natural. One should think of: -
objects as admissible descriptions, - arrows as changes of description,
- orbit invariants as structural content.

Then \(\mathcal{I}(x;T)\) is simply an orbit invariant obtained by
optimizing a nonnegative defect over the orbit.

\subsection{6.3. Derived interpretation}\label{derived-interpretation}

Many classical obstructions arise from failure of a lifting, gluing,
splitting, or commutativity condition. In categorical language, these
appear as: - failure of a diagram to commute, - nonexistence of a global
section, - nontriviality of an extension class, - nonvanishing derived
functor, - nontrivial monodromy, - or obstruction class in a cohomology
group.

The present framework does not replace those structures. Rather, it adds
a new layer: it asks what obstruction remains after minimizing over all
admissible reformulations of the same underlying system.

\section{7. Sheaf-Theoretic and Cocycle
Obstruction}\label{sheaf-theoretic-and-cocycle-obstruction}

This section gives the core local-to-global theorem.

\subsection{7.1. Setup}\label{setup}

Suppose \(x\) admits local descriptions over a cover
\(\mathcal{U} = \{U_i\}\), with transition data \[
g_{ij}
\] defined on overlaps \(U_i \cap U_j\).

Define the triple-overlap defect \[
\kappa_{ijk} := g_{ij} g_{jk} g_{ki}.
\]

In a consistent cocycle, \[
\kappa_{ijk} = e
\] for all triple overlaps. Failure of this identity records a
local-to-global obstruction.

\subsection{7.2. Abstract cocycle obstruction
theorem}\label{abstract-cocycle-obstruction-theorem}

\subsubsection{Theorem 7.1}\label{theorem-7.1}

Assume: 1. The target \(T\) is global trivialization, flattening, or
globally consistent gluing. 2. There exists a nonnegative functional
\(f\) on cocycle defects such that \[
   O(r,x;T) \ge f(\kappa_{ijk})
   \] for every admissible representation \(r\). 3. \(f(\kappa)=0\) if
and only if the cocycle defect vanishes in the admissible class.

If the cocycle class is nontrivial, then \[
\mathcal{I}(x;T) > 0.
\]

\subsubsection{Proof}\label{proof-5}

Suppose for contradiction that \[
\mathcal{I}(x;T)=0.
\] Then for every \(\varepsilon > 0\), there exists
\(r_\varepsilon \in \mathcal{R}(x)\) such that \[
O(r_\varepsilon,x;T) < \varepsilon.
\] By the lower bound hypothesis, \[
f(\kappa_{ijk}(r_\varepsilon)) \le O(r_\varepsilon,x;T) < \varepsilon.
\] Hence \(f(\kappa_{ijk}(r_\varepsilon)) \to 0\). By hypothesis (3),
the cocycle defect must vanish in the admissible limit, contradicting
nontriviality of the cocycle class. Therefore \(\mathcal{I}(x;T) > 0\).
QED.

\subsubsection{Remark}\label{remark-1}

This theorem captures the common pattern that failure of gluing produces
a positive irreducible residue.

\section{8. Gauge-Theoretic
Formulation}\label{gauge-theoretic-formulation}

The gauge-theoretic interpretation is a particularly important special
case.

\subsection{8.1. Local gauge data}\label{local-gauge-data}

Let: - local trivializations be the representations, - gauge
transformations be admissible re-descriptions, - transition functions
encode compatibility, - loop products or cocycles encode global
obstruction.

Then semantic obstruction becomes a gauge-relative obstruction, and
irreducible residue becomes the orbit-minimized gauge-invariant
remainder.

\subsection{8.2. Gauge-invariant residue
theorem}\label{gauge-invariant-residue-theorem}

\subsubsection{Theorem 8.1}\label{theorem-8.1}

Assume the obstruction functional depends only on conjugacy-invariant
data, such as: - holonomy around loops, - Wilson loop residue, - cocycle
defect class, - curvature norms invariant under gauge.

Then the irreducible residue \(\mathcal{I}(x;T)\) is gauge-invariant.

\subsubsection{Proof}\label{proof-6}

Gauge transformations act by conjugation on loop products and cocycle
representatives, but conjugacy-invariant defect functionals do not
change under conjugation. Since the residue is the infimum over the
whole gauge orbit, its value is unchanged by gauge transformation. QED.

\subsection{8.3. Interpretation}\label{interpretation-1}

This theorem captures the abstract skeleton of holonomy-based
obstruction. The same pattern later specializes to normalized invariants
in particular field theories, but the general structure already exists
here.

\section{9. Higher-Categorical Lift}\label{higher-categorical-lift}

The framework admits a natural lift to higher-categorical language.

\subsection{9.1. Representation
infinity-groupoid}\label{representation-infinity-groupoid}

Replace the representation groupoid \(\mathcal{G}\) by an
\(\infty\)-groupoid \(\mathcal{G}_\infty\), where: - 0-morphisms are
representations, - 1-morphisms are admissible re-descriptions, -
2-morphisms are homotopies between re-descriptions, - higher morphisms
encode coherence data.

The orbit structure is then understood homotopy-invariantly.

\subsection{9.2. Higher obstruction
assignments}\label{higher-obstruction-assignments}

An obstruction assignment may now take values in an ordered homotopical
object \(\mathcal{V}\), or be extracted from a spectral invariant,
derived functor, or obstruction tower.

At the simplest level, one may still define \[
\mathcal{I}(x;T) := \inf_{r \in \pi_0(\mathcal{G}_\infty(x))} O(r,x;T),
\] where \(\pi_0\) denotes connected components in the
\(\infty\)-groupoid.

\subsection{9.3. Interpretation}\label{interpretation-2}

The higher-categorical lift is useful when: - admissible transformations
themselves have nontrivial homotopy, - coherence conditions matter, -
obstruction arises through towers of extensions or lifting problems, -
or one wants to compare obstructions up to higher equivalence rather
than strict equality.

The basic residue remains an orbit-minimized obstruction, but now the
orbit is homotopical rather than merely groupoidal.

\section{10. Worked Examples with Full
Proofs}\label{worked-examples-with-full-proofs}

This section gives concrete worked examples illustrating the framework.

\subsection{10.1. Example 1: Gauge holonomy on a circle
bundle}\label{example-1-gauge-holonomy-on-a-circle-bundle}

Let \(M = S^1\). Consider a principal \(U(1)\)-bundle with connection
form \(A\). Let the target \(T\) be global flattening, meaning zero
holonomy around the base circle.

Representations are gauge choices \(A \mapsto A + d\phi\). Let the
obstruction be \[
O(A) := \inf_{n \in \mathbb{Z}} \left| \int_{S^1} A - 2\pi n \right|.
\]

This is the distance of the total holonomy phase from the nearest
trivial \(2\pi \mathbb{Z}\)-phase.

\subsubsection{Claim}\label{claim}

\(O(A)\) is gauge-invariant and equals the irreducible residue.

\subsubsection{Proof}\label{proof-7}

Under gauge transformation, \[
A \mapsto A + d\phi.
\] Hence \[
\int_{S^1} (A + d\phi) = \int_{S^1} A + \int_{S^1} d\phi = \int_{S^1} A,
\] since the integral of an exact form around a closed loop is zero.
Therefore \(O(A)\) is gauge-invariant. Since every gauge-equivalent
representation has the same obstruction value, the infimum over the
orbit equals that common value. Thus \[
\mathcal{I}(A;T) = O(A).
\] If \(\int_{S^1} A \notin 2\pi \mathbb{Z}\), then
\(\mathcal{I}(A;T) > 0\). QED.

\subsubsection{Interpretation}\label{interpretation-3}

This is the simplest instance of a semantic obstruction whose residue is
already fully representation-invariant.

\subsection{10.2. Example 2: Sheaf gluing on a two-set
cover}\label{example-2-sheaf-gluing-on-a-two-set-cover}

Let \(X = U \cup V\), with nonempty overlap \(U \cap V\). Suppose one
has local data \(s_U\) on \(U\) and \(s_V\) on \(V\), and the target
\(T\) is existence of a global section gluing the local pieces.

Let the obstruction be \[
O(r) := \| s_U|_{U \cap V} - s_V|_{U \cap V} \|,
\] for some norm on the overlap data.

\subsubsection{Claim}\label{claim-1}

If the local data differ by a non-removable mismatch on the overlap,
then \(\mathcal{I}(x;T) > 0\).

\subsubsection{Proof}\label{proof-8}

If every admissible re-description preserves the overlap mismatch norm,
then the obstruction functional is constant over the orbit. Hence the
infimum equals that value. If that value is nonzero, the residue is
positive. More generally, if re-description can reduce but not eliminate
the mismatch, then the infimum remains positive. QED.

\subsubsection{Interpretation}\label{interpretation-4}

This example captures the basic logic of gluing defect without invoking
advanced sheaf language.

\subsection{10.3. Example 3: Nonholonomic
control}\label{example-3-nonholonomic-control}

Consider the standard nonholonomic system in \(\mathbb{R}^3\) with
coordinates \((x,y,\theta)\), controlled by \[
\dot{x} = u_1 \cos\theta,\quad
\dot{y} = u_1 \sin\theta,\quad
\dot{\theta} = u_2.
\]

The system cannot move directly in arbitrary directions at each instant.
Let the target \(T\) be arbitrary local reachability.

Define the obstruction \(O(r,x;T)\) to be the instantaneous directional
deficiency in representation \(r\). In any coordinate representation
preserving the control structure, the rank deficiency is nonzero at the
level of immediate motion. Hence the residue is positive: \[
\mathcal{I}(x;T) > 0.
\]

Yet the Lie bracket \[
[X_1,X_2]
\] generates the missing direction, and the system is locally
controllable.

\subsubsection{Claim}\label{claim-2}

This system is irreducible and traversable.

\subsubsection{Proof}\label{proof-9}

The instantaneous rank deficiency cannot be removed by admissible
coordinate change preserving the control structure, so
\(\mathcal{I}(x;T) > 0\). However, Chow-Rashevskii theory implies local
reachability from the bracket-generating property. Thus successful
traversal exists despite positive residue. By Theorem 5.2, the system
lies in the irreducible and traversable regime. QED.

\subsubsection{Interpretation}\label{interpretation-5}

This is a canonical worked example of irreducible obstruction without
fatality.

\subsection{10.4. Example 4: Optimization and convexification
limit}\label{example-4-optimization-and-convexification-limit}

Let \(f : \mathbb{R}^n \to \mathbb{R}\) be a nonconvex objective, and
let \(T\) be global convex solvability.

Suppose admissible representations are diffeomorphic reparameterizations
\(\Phi\), so that in representation \(\Phi\), the objective becomes \[
f_\Phi := f \circ \Phi^{-1}.
\]

Define \[
O(\Phi) := \text{a nonconvexity measure of } f_\Phi,
\] for example the size of the negative part of the Hessian where
defined.

Then \[
\mathcal{I}(f;T) := \inf_\Phi O(\Phi)
\] measures the minimal nonconvexity surviving admissible
reparameterization.

\subsubsection{Claim}\label{claim-3}

If no admissible reparameterization convexifies the objective, then
\(\mathcal{I}(f;T) > 0\).

\subsubsection{Proof}\label{proof-10}

If \(\mathcal{I}(f;T) = 0\), then by definition one could find
admissible reparameterizations along which the nonconvexity measure
tends to zero. Under any closure or attainment hypotheses sufficient for
the problem, that would produce an effectively convex representation,
contradicting the assumption. Therefore the residue is positive. QED.

\subsubsection{Interpretation}\label{interpretation-6}

This example shows how the framework applies even outside topology and
gauge theory.

\section{11. Relation to Prior
Frameworks}\label{relation-to-prior-frameworks}

This section situates the present work relative to well-known
literatures.

\subsection{11.1. Classical obstruction
theory}\label{classical-obstruction-theory}

Classical obstruction theory studies whether one may extend or section a
partial construction, with obstructions often living in cohomology
groups. The present framework does not replace obstruction theory.
Instead, it adds a representational layer: one asks not only whether an
obstruction exists, but whether it can be diminished or removed by
admissible reformulation. The irreducible residue is the orbit-minimized
shadow of classical obstruction.

\subsection{11.2. Sheaf theory and
descent}\label{sheaf-theory-and-descent}

The gluing language in Section 7 aligns naturally with sheaf-theoretic
descent. Čech cocycles, nontrivial cohomology classes, and descent
obstructions are concrete sources of positive residue when the target is
global consistency or trivialization.

\subsection{11.3. Gauge theory}\label{gauge-theory}

Gauge theory already organizes local descriptions into orbits under
gauge equivalence. Holonomy, curvature, cocycle data, and topological
charge supply natural obstruction functionals. The current framework
extracts the general principle that the meaningful obstruction is the
orbit-minimized gauge-invariant residue.

\subsection{11.4. Control theory}\label{control-theory}

Nonholonomic systems provide a paradigm case in which positive residue
does not imply impossibility. The irreducible and traversable regime is
standard in control practice but seldom elevated to a general structural
principle.

\subsection{11.5. Category theory and higher
structures}\label{category-theory-and-higher-structures}

The representation groupoid and its higher-categorical lift place the
framework in a natural categorical setting. The irreducible residue can
be understood as an orbit-invariant lower envelope across a category of
admissible descriptions.

\section{12. Design Principle and
Interpretation}\label{design-principle-and-interpretation}

The framework suggests the following general principle.

\subsubsection{Principle of Semantic
Irreducibility}\label{principle-of-semantic-irreducibility}

The meaningful obstruction in a problem is not \(O(r,x;T)\) for one
chosen representation \(r\), but the orbit-minimized quantity \[
\mathcal{I}(x;T).
\]

\subsubsection{Operational
reformulation}\label{operational-reformulation}

If \(\mathcal{I}(x;T)=0\), the obstruction is removable and one should
search for a better representation.

If \(\mathcal{I}(x;T)>0\), then the question changes. One must determine
whether the problem is fatal or traversable.

This turns many design questions from: Can the obstruction be removed?

into: Does the orbit-minimized residue vanish, and if not, can traversal
be organized around it?

\section{13. Conclusion}\label{conclusion}

We have introduced a general framework for semantic obstruction and
irreducible residue. The framework formalizes: - representation-relative
obstruction, - orbit-minimized structural remainder, - canonical
decomposition, - orbit invariance, - existence under mild hypotheses, -
local-to-global obstruction via cocycle defect, - and a trichotomy
separating removable, fatal, and traversable regimes.

The framework is intentionally general. It does not solve any
domain-specific open problem by itself. Instead, it provides a reusable
front-end for formulating such problems in a way that cleanly separates
representational artifact from structural content. Its central
contribution is the rigorous formalization of the irreducible but
traversable regime.

\section{References}\label{references}

A concise submission-ready bibliography can begin as follows.

Baez, J. C., and Schreiber, U. Higher gauge theory. Categories in
algebra, geometry and mathematical physics.

Bott, R., and Tu, L. W. Differential Forms in Algebraic Topology.

Grothendieck, A. Séminaire de Géométrie Algébrique.

Husemoller, D. Fibre Bundles.

Lurie, J. Higher Topos Theory.

Mac Lane, S. Categories for the Working Mathematician.

May, J. P. A Concise Course in Algebraic Topology.

Montgomery, R. A Tour of Subriemannian Geometries, Their Geodesics and
Applications.

Steenrod, N. The Topology of Fibre Bundles.

Whitehead, G. W. Elements of Homotopy Theory.

Arora, S., and Barak, B. Computational Complexity: A Modern Approach.

Lee, J. M. Introduction to Smooth Manifolds.

Kashiwara, M., and Schapira, P. Categories and Sheaves.

\section{Appendix A. Universal Front-End for Clay Problems and Related
Programs}\label{appendix-a.-universal-front-end-for-clay-problems-and-related-programs}

This appendix explains how the present framework can be used as a
universal front-end for broad research programs without claiming that
the abstract framework itself solves any one of them.

\subsection{A.1. General front-end
recipe}\label{a.1.-general-front-end-recipe}

To apply the framework to a problem family, specify five ingredients:

\begin{enumerate}
\def\labelenumi{\arabic{enumi}.}
\tightlist
\item
  System space \(X\)
\item
  Target \(T\)
\item
  Representation class \(\mathcal{R}(x)\) and admissible groupoid
  \(\mathcal{G}\)
\item
  Obstruction functional \(O\)
\item
  Success predicate \(P\)
\end{enumerate}

Once these are fixed, compute or estimate: \[
\mathcal{I}(x;T) = \inf_{r \in \mathcal{R}(x)} O(r,x;T).
\]

Then ask:

\begin{enumerate}
\def\labelenumi{\arabic{enumi}.}
\tightlist
\item
  Does \(\mathcal{I}(x;T)=0\)?
\item
  If not, is the pair fatal or traversable?
\item
  Under scale change or coarse-graining, does the residue persist,
  dilute, or amplify?
\end{enumerate}

This gives a disciplined front-end before entering domain-specific
technical machinery.

\subsection{A.2. Yang-Mills-type
program}\label{a.2.-yang-mills-type-program}

A front-end formulation may be organized as follows.

System: local gauge data, holonomy data, or quantum field
configurations.

Target: global flattening, mass-gap-compatible spectral control, or
elimination of trivialization failure.

Representations: local trivializations, gauge choices, coarse-grained
encodings, representation-theoretic descriptions.

Obstruction: loop holonomy defect, cocycle defect, Wilson loop residue,
normalized character defect, or mean-zero contraction defect.

Residue: orbit-minimized nontrivial holonomy or normalized obstruction.

Front-end question: Does a nonzero normalized residue persist under
admissible scale passage?

This does not prove the mass gap. What it does is stabilize the
conceptual target and prevent confusion between raw small-scale defect
and normalized structural content.

\subsection{A.3. Navier-Stokes-type
program}\label{a.3.-navier-stokes-type-program}

System: velocity fields, vorticity fields, or scale-localized fluid
states.

Target: global regularity, suppression of blow-up, or elimination of
unresolved singular concentration.

Representations: coordinate systems, scale decompositions, wavelet
frames, geometric flow-adapted descriptions.

Obstruction: concentration defect, vortex stretching imbalance,
non-removable cascade concentration, or singularity indicator.

Residue: minimal concentration obstruction surviving admissible
reformulation.

Front-end question: Are observed blow-up signals representational
artifacts, or is there an orbit-invariant concentration residue?

This does not settle regularity. It clarifies what kind of object one
would need to show vanishes, remains bounded, or becomes traversable
under a suitable reformulation.

\subsection{A.4. Riemann-Hypothesis-type
program}\label{a.4.-riemann-hypothesis-type-program}

System: explicit formula data, zero distributions, spectral traces,
phase-encoded arithmetic objects.

Target: critical symmetry, elimination of off-line spectral defect, or
restoration of a canonical balanced description.

Representations: spectral formulations, transformed explicit formulas,
geometric encodings, dynamical systems analogues, phase-unwrapped
descriptions.

Obstruction: critical-line deviation functional, asymmetry residue,
phase mismatch, or non-removable spectral imbalance.

Residue: the minimal spectral asymmetry surviving admissible
reformulation.

Front-end question: Does the orbit-minimized spectral defect vanish?

This does not prove the hypothesis. It provides a formal way to
distinguish representational asymmetry from intrinsic asymmetry.

\subsection{A.5. P versus NP-type
program}\label{a.5.-p-versus-np-type-program}

System: problem families, proof systems, encoding models, circuit
classes.

Target: tractable solution class, equivalence collapse, or elimination
of computational hardness witness.

Representations: different encodings, reductions, proof frameworks,
circuit bases, algebraic lifts.

Obstruction: lower-bound witness, barrier phenomenon, nonuniform
hardness defect, compression obstruction.

Residue: the minimal hardness obstruction surviving admissible
representation change.

Front-end question: Are currently observed lower-bound barriers merely
framework-dependent, or do they contain a positive irreducible residue
across admissible encodings?

Again, this does not solve the problem. It gives a clean language for
distinguishing representational barriers from invariant hardness.

\subsection{A.6. General lessons from the front-end
viewpoint}\label{a.6.-general-lessons-from-the-front-end-viewpoint}

Across all these programs, the same workflow emerges.

Step 1: Define the target before choosing the representation.

Step 2: Describe the admissible representation orbit.

Step 3: Choose an obstruction functional that reflects the actual
difficulty, not just a convenient proxy.

Step 4: Compute or estimate the irreducible residue.

Step 5: Decide whether the regime is removable, fatal, or traversable.

Step 6: Only then apply specialized analytic, algebraic, topological, or
computational machinery.

This helps prevent the common failure mode in which one chases an
obstruction that is merely representational while ignoring the
orbit-invariant remainder.

\subsection{A.7. Why this front-end is
useful}\label{a.7.-why-this-front-end-is-useful}

The front-end is useful because it does three things at once.

First, it forces one to say explicitly what counts as an admissible
reformulation.

Second, it separates the target from the current descriptive language.

Third, it converts vague claims like ``this seems fundamental'' or
``this is probably an artifact'' into a well-defined mathematical
question about whether \(\mathcal{I}(x;T)\) vanishes.

\subsection{A.8. Meta-principle}\label{a.8.-meta-principle}

The universal lesson is:

The right question is not whether an obstruction appears in one
formulation, but whether its orbit-minimized residue survives all
admissible formulations.

That sentence can function as the conceptual front-end to broad classes
of difficult problems.

\section{Appendix B. Submission Polishing
Notes}\label{appendix-b.-submission-polishing-notes}

\subsection{B.1. Suggested titles}\label{b.1.-suggested-titles}

Here are several title variants, depending on tone.

Formal: Semantic Obstruction and Irreducible Residue: A
Category-Theoretic Framework for Representation-Invariant Constraints

More conceptual: Irreducible Residue: A General Theory of
Representation-Invariant Obstruction and Traversability

More applied: Semantic Obstruction, Traversability, and Local-to-Global
Failure: A Unifying Framework

\subsection{B.2. Suggested short
title}\label{b.2.-suggested-short-title}

Semantic Obstruction and Irreducible Residue

\subsection{B.3. Suggested opening
footnote}\label{b.3.-suggested-opening-footnote}

This paper develops a general formal framework. It does not claim to
solve any particular open problem by abstract argument alone; rather, it
isolates a reusable invariant structure intended to organize
domain-specific proof programs.

\subsection{B.4. Suggested arXiv
categories}\label{b.4.-suggested-arxiv-categories}

Primary candidates: - math.CT - math.AT - math-ph

Possible secondary cross-listing depending on final examples: - math.OC
- cs.CC

\subsection{B.5. Suggested author
note}\label{b.5.-suggested-author-note}

If you want a philosophical but disciplined positioning sentence:

This work is foundational rather than problem-specific. Its aim is to
formalize a general mathematical structure underlying a wide range of
local-to-global, gauge-relative, and representation-sensitive
obstructions.

\section{Appendix C. Copy-Paste Preservation
Notes}\label{appendix-c.-copy-paste-preservation-notes}

To preserve the mathematics when copying into plain text or Markdown
environments:

\begin{enumerate}
\def\labelenumi{\arabic{enumi}.}
\tightlist
\item
  Keep all formulas either inline as \texttt{\$...\$} or display as
  \texttt{\$\$...\$\$}.
\item
  Avoid nonstandard theorem environments if moving across editors.
\item
  Keep equation punctuation outside the LaTeX delimiters when possible.
\item
  Use named definitions in plain text headers rather than custom macros
  unless you control the target renderer.
\end{enumerate}

Example safe style:

Definition 3.1. The irreducible residue is \[
\mathcal{I}(x;T) := \inf_{r \in \mathcal{R}(x)} O(r,x;T).
\]

This style will survive most copy-paste pipelines.

\end{document}
